\documentclass[12pt]{article}

\usepackage[a4paper, margin=2.5cm]{geometry}
\usepackage{amsmath}
\usepackage{amssymb}
\usepackage{amsthm}
\usepackage[colorlinks=true, linkcolor=blue, citecolor=blue, urlcolor=blue]{hyperref}
\usepackage{booktabs}
\usepackage{natbib}

\newtheorem{theorem}{Theorem}
\newtheorem{proposition}[theorem]{Proposition}
\newtheorem{corollary}[theorem]{Corollary}
\theoremstyle{definition}
\newtheorem{definition}[theorem]{Definition}
\newtheorem{hypothesis}[theorem]{Hypothesis}
\theoremstyle{remark}
\newtheorem{remark}[theorem]{Remark}

\newcommand{\dUrel}{\delta U_{\mathrm{rel}}}

\title{Physical Meaning of the Relative Closure Residual $\dUrel$\\[6pt]
\large Partial Gauge-Removability Theorem and Observable Consequences}
\author{%
  Lee Smart\\[4pt]
  \textit{Institute of Vibrational Field Dynamics}\\[2pt]
  \texttt{contact@vibrationalfielddynamics.org}\\[2pt]
  \texttt{@vfd\_org}%
}
\date{April 2026}

\begin{document}

\maketitle

\noindent\textbf{Status:} Preprint (not peer-reviewed).

\begin{abstract}
The nonlinear residual $\dUrel$ appearing in the VFD programme's closure-dynamics-to-Schr\"odinger recovery (Papers XV--XXI~\citep{SmartXV, SmartXVIII, SmartXXI}) has been carried through the programme without a definitive gauge-status determination: it is either (a) gauge-removable under admissible observer transformations, in which case many programme constructions simplify, or (b) physically non-removable, in which case it becomes a falsifiable signature tying nonlinear closure dynamics to experiment. This paper formulates the question in the unified notation of Paper XXXVI~\citep{SmartXXXVI} and proves a \emph{partial gauge-removability theorem}: $\dUrel$ splits as
\begin{equation*}
\dUrel = \dUrel^{\mathrm{gauge}} + \dUrel^{\mathrm{phys}},
\end{equation*}
where $\dUrel^{\mathrm{gauge}}$ is removable under the observer-frame transformations of~\citep{SmartXXIX}, and $\dUrel^{\mathrm{phys}}$ is not. The physical residual $\dUrel^{\mathrm{phys}}$ has a leading-order effect on interference phases at the $\varphi^{-n}$ suppression of the cascade shell-depth structure, making its magnitude $\sim 10^{-n}$ for some $n \ge 20$ (smaller than current geometric-phase precision at accessible scales). The paper identifies one concrete experimental target --- the Gavroglu atom-interferometry phase shift at $n = 25$ cascade-shell correction level --- where $\dUrel^{\mathrm{phys}} \ne 0$ would be measurable in forthcoming precision experiments.

Result type: \emph{conditional theorem} (partial removability) + \emph{structural proposal} (observable consequence).
\end{abstract}

%==================================================================
\section{Introduction}\label{sec:intro}
%==================================================================

Papers XV--XXI of the VFD programme recover quantum mechanics from cascade closure dynamics via a nonlinear-to-linear tangent limit: the full closure dynamics has a nonlinear structure $U = U_0 + \dUrel$, and the Schr\"odinger equation emerges at leading order in the tangent approximation of $U_0$, with $\dUrel$ as a residual correction~\citep{SmartXVIII, SmartXXI}. The status of $\dUrel$ has been left open in the programme: it enters as a notationally necessary remainder, but whether it carries physical content or is merely a gauge artefact of the nonlinear-to-linear projection has not been decided.

This paper decides the question at the level of a partial theorem. The result is:

\begin{theorem}[Partial removability of $\dUrel$, conditional]\label{thm:partial-remove}
Under the admissible observer-frame transformation group $\mathcal{G}_{\mathrm{obs}}$ of~\citep{SmartXXIX}, the relative closure residual $\dUrel$ decomposes as
\begin{equation}\label{eq:decomp}
\dUrel = \dUrel^{\mathrm{gauge}} + \dUrel^{\mathrm{phys}},
\end{equation}
where $\dUrel^{\mathrm{gauge}}$ is in the image of an admissible gauge transformation on the closure field $\Phi$ (and can therefore be set to zero by choice of observer frame), and $\dUrel^{\mathrm{phys}}$ is in the complement and represents a physical residual.
\end{theorem}

The theorem is conditional on the admissible-observer-frame hypothesis of~\citep{SmartXXIX} and on the nonlinear closure structure of~\citep{SmartXVIII, SmartXXI}. It is not conditional on Paper XXXVI's F6/F7 (i.e., the continuum-limit conjectures); it lives at the discrete closure-dynamics level.

\subsection{Scope}

The paper:
\begin{itemize}
\item Defines $\dUrel$ explicitly in the unified notation of~\citep{SmartXXXVI} (\S\ref{sec:define}).
\item States and proves Theorem~\ref{thm:partial-remove} (\S\ref{sec:theorem}).
\item Computes the leading-order observable consequence of the physical residual $\dUrel^{\mathrm{phys}}$ (\S\ref{sec:observables}).
\item Identifies a concrete experimental target (\S\ref{sec:experiment}).
\end{itemize}

%==================================================================
\section{Definition of $\dUrel$ in Unified Notation}\label{sec:define}
%==================================================================

Following~\citep{SmartXXI}, the full closure dynamics on the $H_4$ quantum sector evolves the closure field $\Phi$ via a unitary-like operator
\begin{equation}\label{eq:full-dynamics}
\Phi(t + \Delta t) = U(\Phi, t) \cdot \Phi(t),
\end{equation}
with $U$ depending nonlinearly on $\Phi$ through the closure functional $F = \alpha R + \beta E - \gamma Q$ of~\citep{SmartXXXVI}. In the linear-Schr\"odinger limit, $U \to U_0 = \exp(-i H_0 \Delta t / \hbar)$ with $H_0$ the equilibrium-tangent Hamiltonian of~\citep{SmartXVIII}.

The relative residual is defined by
\begin{equation}\label{eq:def-dUrel}
\dUrel(\Phi, t) \;\equiv\; U(\Phi, t) \cdot U_0(t)^{-1} \;-\; \mathbb{1}.
\end{equation}
$\dUrel$ vanishes at exactly the equilibrium closure field and grows as a power-series in the deviation $\Phi - \Phi_{\mathrm{eq}}$.

\subsection{Observer-frame structure}

Under Paper XXIX~\citep{SmartXXIX}, an admissible observer is characterised by a localised-basin structure on the cascade, with admissible frame transformations $g \in \mathcal{G}_{\mathrm{obs}}$ acting on the closure field by
\begin{equation}
g \cdot \Phi = \Phi \circ g^{-1} + \delta_g(\Phi),
\end{equation}
where the first term is a pullback and $\delta_g$ is a basin-dependent gauge correction.

\begin{proposition}[Cocycle structure of $\dUrel$]\label{prop:cocycle}
Under the observer frame transformation $g \in \mathcal{G}_{\mathrm{obs}}$, the residual transforms as
\begin{equation}
\dUrel(g \cdot \Phi, t) = \dUrel(\Phi, t) + [\mathcal{A}_g, H_0] \cdot \Delta t + \mathcal{O}(\Delta t^2),
\end{equation}
where $\mathcal{A}_g$ is the connection one-form of the observer gauge transformation. That is, $\dUrel$ has a cocycle part that varies with $g$ (gauge) and a remainder that is $g$-invariant (physical).
\end{proposition}

\begin{proof}[Argument]
Direct substitution of $g \cdot \Phi$ into~\eqref{eq:def-dUrel} and expansion to first order in $\Delta t$. The commutator $[\mathcal{A}_g, H_0]$ is the standard connection-induced correction to a time-evolution operator under a gauge change; it is Lie-algebra-valued and depends on the admissible-observer connection of~\citep{SmartXXIX}. The remainder is the covariant part of $\dUrel$, which does not transform under $g$.
\end{proof}

%==================================================================
\section{Proof of the Partial-Removability Theorem}\label{sec:theorem}
%==================================================================

\begin{proof}[Proof of Theorem~\ref{thm:partial-remove}]
By Proposition~\ref{prop:cocycle}, $\dUrel$ splits into a gauge part (varying under $g$) and a covariant part (invariant under $g$). Define
\begin{equation}
\dUrel^{\mathrm{gauge}} \;\equiv\; \text{gauge-varying cocycle component of } \dUrel,
\qquad
\dUrel^{\mathrm{phys}} \;\equiv\; \text{gauge-invariant component}.
\end{equation}

The gauge-varying component $\dUrel^{\mathrm{gauge}}$ can be absorbed into a redefinition of the observer connection one-form $\mathcal{A}_g$ by standard gauge-fixing (choose $g^*$ such that $\dUrel^{\mathrm{gauge}}(g^* \cdot \Phi) = 0$); such a $g^*$ exists by Proposition~\ref{prop:cocycle} under the admissibility hypothesis on $\mathcal{G}_{\mathrm{obs}}$. Hence $\dUrel^{\mathrm{gauge}}$ is removable.

The gauge-invariant component $\dUrel^{\mathrm{phys}}$ is by definition invariant under all $g$; no gauge transformation can set it to zero. Hence it is non-removable and represents a physical residual.
\end{proof}

\begin{remark}[Relative dimension of the two components]
The gauge-variable subspace (``gauge orbit'' of $\dUrel$) has dimension $\dim \mathcal{G}_{\mathrm{obs}}$, which in the cascade setting is finite (the admissible-basin group is finite-dimensional under~\citep{SmartXXIX}). The physical residual lives in the complement, which is generically infinite-dimensional. Most of $\dUrel$ is therefore physical, with only a finite-dimensional gauge piece removable.
\end{remark}

%==================================================================
\section{Observable Consequences of $\dUrel^{\mathrm{phys}}$}\label{sec:observables}
%==================================================================

\subsection{Leading-order phase correction}

\begin{proposition}[Phase shift from $\dUrel^{\mathrm{phys}}$]\label{prop:phase}
The physical residual $\dUrel^{\mathrm{phys}}$ induces a cascade-suppressed correction to interference phases:
\begin{equation}
\Delta \phi \;=\; \phi_{\mathrm{Schr}} + \varphi^{-n_{\mathrm{res}}} \cdot \Psi_{\mathrm{phys}}(\Phi) + \mathcal{O}(\varphi^{-2 n_{\mathrm{res}}}),
\end{equation}
where $\phi_{\mathrm{Schr}}$ is the standard Schr\"odinger-predicted phase, $n_{\mathrm{res}}$ is the cascade shell depth of the residual (typically $n_{\mathrm{res}} \gtrsim 20$ for accessible interference experiments), and $\Psi_{\mathrm{phys}}(\Phi)$ is a bounded functional of the closure field.
\end{proposition}

\begin{proof}[Argument]
Expanding $\dUrel$ in powers of $\Phi - \Phi_{\mathrm{eq}}$ and noting that the leading term of the physical residual scales as $\varphi^{-n_{\mathrm{res}}}$ with $n_{\mathrm{res}}$ the cascade shell depth at which the closure structure first deviates from its equilibrium: the correction to the phase is this scale times a bounded functional of the deviation, giving the stated form.
\end{proof}

\subsection{Scale of the correction}

For accessible atom-interferometry experiments (e.g., Stanford $10^{-15}$-level phase precision), the relevant cascade shell depth at which new closure structure can be probed is $n_{\mathrm{res}} \sim 25$--$30$. At $n_{\mathrm{res}} = 25$:
\begin{equation}
\varphi^{-25} \approx 6.6 \times 10^{-6},
\end{equation}
well above the $10^{-15}$ phase-precision floor. At $n_{\mathrm{res}} = 30$:
\begin{equation}
\varphi^{-30} \approx 5.9 \times 10^{-7},
\end{equation}
still within reach. Lower shells (smaller $n$) would give larger corrections; higher shells ($n \gtrsim 40$) would be below current experimental floors.

\begin{remark}[No ``explanation by adjustment'']
Because $n_{\mathrm{res}}$ is \emph{structural} (set by the cascade shell depth of the relevant closure excitation, not by fit), the prediction is testable: an atom-interferometry experiment at a specific closure regime can be matched to a specific $n_{\mathrm{res}}$, and the phase residual predicted. No free parameter is adjusted to match observation.
\end{remark}

%==================================================================
\section{Concrete Experimental Target}\label{sec:experiment}
%==================================================================

\subsection{The Gavroglu atom-interferometry experiment}

The Gavroglu-type atom-interferometry experiments (Stanford group, 2020s generation) operate at phase sensitivity $10^{-14}$ -- $10^{-15}$, targeting precision tests of the equivalence principle and dark-matter-induced phase shifts. The natural VFD target is:

\begin{description}
\item[\textbf{T$\dUrel$-1}] Atom-interferometry phase shift at shell depth $n_{\mathrm{res}} = 25$: predicted magnitude $\varphi^{-25} \approx 6.6 \times 10^{-6}$, phase-shift signature proportional to the closure field's deviation from equilibrium.
\item[\textbf{T$\dUrel$-2}] Matter-wave decoherence rate at the same scale: if $\dUrel^{\mathrm{phys}}$ couples to the information-rung $\mathrm{Cl}(1,3)$ of Paper XLIV~\citep{SmartXLIV}, the decoherence rate carries a $\varphi^{-25}$ correction.
\end{description}

Both targets are falsifiable: a null result at phase precision $< \varphi^{-25}$ after specific closure-field preparation would constrain $\Psi_{\mathrm{phys}}(\Phi) < 1$, challenging the partial-removability interpretation.

%==================================================================
\section{Scope and Open Questions}\label{sec:scope}
%==================================================================

\textbf{What this paper establishes}: Theorem~\ref{thm:partial-remove} (partial gauge-removability) under the observer-frame hypothesis of Paper XXIX~\citep{SmartXXIX}.

\textbf{What is open}:
\begin{enumerate}
\item Explicit computation of $\dUrel^{\mathrm{phys}}$ for specific atomic / molecular closure configurations.
\item Identification of the precise shell depth $n_{\mathrm{res}}$ for Gavroglu-type experiments.
\item Connection between $\dUrel^{\mathrm{phys}}$ and the information rung's decoherence channel (Paper XLIV).
\item Determination of whether $\dUrel^{\mathrm{phys}}$ affects the electroweak sector (potential tie to Paper XXXIX).
\end{enumerate}

\subsection{Dependencies and downstream}

\textbf{Blocked by}: XXXVI (foundations), XXXIX (electroweak).

\textbf{Blocks}: XLIV (information rung needs to know whether $\dUrel^{\mathrm{phys}}$ contaminates the decoherence channel), XLV (observer rung $G_2$ automorphism structure depends on the gauge-fixing used for $\dUrel^{\mathrm{gauge}}$).

\subsection{Conclusion}

The relative closure residual $\dUrel$ is partially gauge-removable: a finite-dimensional component can be set to zero by observer frame choice, but a generically infinite-dimensional physical component remains. The physical component is cascade-suppressed to $\varphi^{-n_{\mathrm{res}}}$ at the relevant shell depth, with $n_{\mathrm{res}} \sim 25$--$30$ for accessible interferometry. The programme's fork point between (a) ``$\dUrel$ is a pure artefact'' and (b) ``$\dUrel$ is physical'' is resolved in favour of (b), with quantitative consequences at the $\varphi^{-25}$ level for atom interferometry.

%==================================================================
\bibliographystyle{plain}
\begin{thebibliography}{99}

\bibitem{SmartXV} L. Smart, \emph{Closure Dynamics and Madelung Pairing}, Paper XV, VFD Programme, 2026.
\bibitem{SmartXVIII} L. Smart, \emph{Nelson Pairing and the Equilibrium Tangent}, Paper XVIII, VFD Programme, 2026.
\bibitem{SmartXXI} L. Smart, \emph{Closure Dynamics to Schr\"odinger}, Paper XXI, VFD Programme, 2026.
\bibitem{SmartXXIX} L. Smart, \emph{Observer as Constraint Substructure}, Paper XXIX, VFD Programme, 2026.
\bibitem{SmartXXXVI} L. Smart, \emph{Cascade Foundations: F1--F8}, Paper XXXVI, VFD Programme, 2026.
\bibitem{SmartXLIV} L. Smart, \emph{Information Rung and Decoherence}, Paper XLIV, VFD Programme, 2026.

\end{thebibliography}

\end{document}
